% ============================================================
%  related.tex — 相关工作:批判性归类(引 ② taxonomy),不罗列堆砌
%  纪律:每条引用来自打开过原文的语料清单;说清本工作与各路线的差异。
% ============================================================

\section{Related Work}
\label{sec:related}

% 按 ② SurveyWeaver 的 method taxonomy 分小节,每节末点出「与本工作的差异」。

\subsection{Route A: <名称, 如 Imitation Learning for Manipulation>}
Representative works~\cite{<key>,<key>} \ldots. Their strength is \ldots,
but they \textbf{differ from ours in} \ldots.

\subsection{Route B: <名称, 如 Tactile / Multi-modal Sensing>}
\ldots~\cite{<key>}. In contrast, we \ldots

\subsection{Route C: <名称, 如 Language-conditioned / VLA Policies>}
\ldots~\cite{<key>}. Unlike these, our approach \ldots

% [定位小结] 一段话把本工作放进版图,明确差异化(回应 ① 的 gap 陈述)。
\paragraph{Positioning.} Compared with the above, MethodName is distinct in that
\ldots, which directly addresses the gap identified in Sec.~\ref{sec:intro}.
